\subsection{Практические задания: реакция на события}
Этот блок выполняется после раздела про \verb|callback(event)|. Здесь важно не просто менять цвет светодиодов, а корректно реагировать на события автопилота и не использовать блокирующие задержки.

\begin{taskbox}[Реакция на событие посадки]
\indent Требуется реализовать обработчик событий, который при получении события \verb|Ev.COPTER_LANDED| полностью отключает светодиодную индикацию. Реализация не должна использовать блокирующие задержки.

\noindent\textbf{Ожидаемый результат:} после посадки все светодиоды выключены.
\end{taskbox}

\begin{taskbox}[Аварийная индикация при низком напряжении]
\indent Необходимо реализовать обработку события \verb|Ev.LOW_VOLTAGE2|, при котором все активные таймеры останавливаются, а светодиодная линейка переводится в устойчивое состояние красной индикации.

\noindent\textbf{Ожидаемый результат:} линейка постоянно светится красным цветом, сигнализируя об аварийном режиме.
\end{taskbox}
